

\documentclass[letterpaper,12pt]{article}

\RequirePackage{amsthm,amsmath,amssymb,amscd}

\usepackage{latexsym}
\usepackage[empty]{fullpage}
\usepackage{titlesec}
\usepackage{marvosym}
\usepackage[usenames,dvipsnames]{color}
\usepackage{verbatim}
\usepackage{enumitem}
\usepackage[pdftex]{hyperref}
\usepackage{fancyhdr}
\usepackage{multicol}
\usepackage{lipsum}

\titlespacing*{\section}{0pt}{.4\baselineskip}{\baselineskip}

\hypersetup{colorlinks=true, citecolor=blue, urlcolor=blue,  linkcolor=blue, linktocpage=true}


\pagestyle{fancy}
\fancyhf{} % clear all header and footer fields
\fancyfoot{}
\renewcommand{\headrulewidth}{0pt}
\renewcommand{\footrulewidth}{0pt}

\usepackage[margin=.7in]{geometry}

% Adjust margins
%\addtolength{\oddsidemargin}{-0.375in}
%\addtolength{\evensidemargin}{-0.375in}
%\addtolength{\textwidth}{.5in}
%\addtolength{\topmargin}{-.2in}
%\addtolength{\textheight}{1in}

\urlstyle{same}

\raggedbottom
\raggedright
\setlength{\tabcolsep}{0in}


% Sections formatting
\titleformat{\section}
{
  \vspace{0pt}\scshape\raggedright\large
} {}{0em}{} [\color{black}\titlerule \vspace{0pt}]

%-------------------------
% Custom commands
\newcommand{\resumeItem}[2]{
  \item\small{
    \textbf{#1}{: #2 \vspace{-2pt}}
  }
}

\newcommand{\resumeSubheading}[4]{
  \vspace{-1pt}\item
    \begin{tabular*}{0.97\textwidth}{l@{\extracolsep{\fill}}r}
      \textbf{#1} & #2 \\
      \textit{\small#3} & \textit{\small #4} \\
    \end{tabular*}\vspace{-pt}
}

\newcommand{\resumeSubItem}[2]{\resumeItem{#1}{#2}\vspace{-4pt}}

\renewcommand{\labelitemii}{$\circ$}

\newcommand{\resumeSubHeadingListStart}{\begin{itemize}[leftmargin=*]}
\newcommand{\resumeSubHeadingListEnd}{\end{itemize}}
\newcommand{\resumeItemListStart}{\begin{itemize}}
\newcommand{\resumeItemListEnd}{\end{itemize}\vspace{-5pt}}

%-------------------------------------------
%%%%%%  CV STARTS HERE  %%%%%%%%%%%%%%%%%%%%%%%%%%%%

\setlength{\columnsep}{-.1cm}
\begin{document}

%----------HEADING-----------------
\begin{tabular*}{\textwidth}{l@{\extracolsep{\fill}}r}
  \textbf{\Huge Nicholas Eisenberg}& Email : \href{mailto:nickeisenberg@gmail.com}{nickeisenberg@gmail.com}\\
  2945 Mondavi Ct, Las Vegas, Nevada 89117 & Mobile : 702-354-0866 \\
  & \href{https://www.linkedin.com/in/nick-eisenberg/}{Linkedin} \\
  & \href{https://github.com/nickeisenberg}{GitHub}
\end{tabular*}


%------exams------




%-----------EDUCATION-----------------
%\section{Summary}
%I am a PhD student at Auburn University studying Probability and Stochastic Partial Differential Equations under the supervision of \href{http://webhome.auburn.edu/~lzc0090/}{Dr. Le Chen}. I have one publication in the journal {\em Stochastics and Partial Differential Equations: Analysis and Computations}, and I also have another paper with proven results that is now being formatted for submission. I am proficient  with Python which was used extensively throughout my  PhD studies. Although  my work experience has primarily revolved around research, my goal is to achieve a career in the field of data science, and recently I have begun self-study with machine learning. I am looking forward to expanding my knowledge and collaborating with others along the way. 
\section{Summary}
I am a PhD student at Auburn University studying Probability and Stochastic Partial Differential Equations under the supervision of \href{http://webhome.auburn.edu/~lzc0090/}{Dr. Le Chen}. I have one publication in the journal {\em Stochastics and Partial Differential Equations: Analysis and Computations}, and I also have another paper with proven results that is now being formatted for submission. I am proficient in Excel and in Python which was used extensively throughout my  PhD studies and I have also begun learning deep learning through self study.
\section{Education}


 \begin{itemize}
   \item  \noindent Ph.D., \textbf{Auburn University} \hfill Thesis Defense: July 2022 \\
   Field of Study: Stochastic Partial Differential Equtions  \\ 
   Doctoral Adivsor: \href{http://webhome.auburn.edu/~lzc0090/}{Dr. Le Chen}
   
   \item \noindent M.S., \textbf{University of Nevada, Las Vegas } \hfill Jan 2015 - Aug 2018 \\
   		Major: Applied Mathematics \\ 
   		GPA: 3.9
   		
   	\item \noindent B.S., \textbf{University of Nevada, Las Vegas} \hfill Aug 2011 - Dec 2015 \\
   	    	Major: Applied Mathematics \\ 
   	     	GPA: 3.69 \\

     \end{itemize}
	
%\begin{multicols}{3}
%	\begin{itemize}
%  \item \noindent M.S., \textbf{University of Nevada, Las Vegas } \\
%		Major: Applied Mathematics \\ 
%		GPA: 3.9 \\
%		Jan 2015 - Aug 2018 
%      
%    \item \noindent B.S., \textbf{University of Nevada, Las Vegas} \\
%    	Major: Applied Mathematics \\ 
%     	GPA: 3.69 \\
%     	Aug 2011 - Dec 2015 \\ \vspace{1in}
%      
%         \item \textbf{ Actuarial Exams} \\
%   			- P, FM and MFE 
%  \end{itemize}
%\end{multicols}   

%\begin{itemize} 
%%     \item \textbf{ Actuarial Exams} \\
%%     Have passed Exam P, Exam FM and Exam MFE 
%\end{itemize} 

 \section{Publications and Current Projects}

\begin{itemize}[]
	\item \noindent N. Eisenberg and L. Chen. \href{https://link.springer.com/epdf/10.1007/s40072-022-00258-6?sharing_token=615BToOAX2__hha95bBy5Pe4RwlQNchNByi7wbcMAY7h99-9CUgKCJ_P7MgI29SutxQPHeHVRWjEYXbYxT-ZuTjdZ_sthIDAaYOeQZxG9dXLFbji62022XQenyFaNn3mbf7oxkSN8vMrt7ARNnKURUK0-EU0s4A4ilj_xTf2X_8\%3D}{Interpolating the Stochastic Heat and Wave Equations with Time-independent Noise:
		Solvability and Exact Asymptotics} - {\em Stochastics and Partial Differential Equations: Analysis and Computations}
	
	\item \noindent N. Eisenberg and L. Chen. Invariant measures for the stochastic heat equation on $\mathbb{R}^d$ without drift term. Formatting for submission.
	
	\item \noindent N. Eisenberg and L. Chen. The interpolated stochastic heat and wave equation with super-linear drift term: solvability, regualrity and moment bounds. In preparation.
	
\end{itemize}

\section{Skills}
%\begin{itemize}[itemsep = -2pt, topsep=0pt]
\begin{itemize}[]
	\item \textbf{Python, Mathematica, Excel, R, GitHub and LaTex} \\
	\vspace{.05in} 
	- Proficient with the Python packages NumPy, math, pandas and matpoltlib.pyplot. 
	
	\item \textbf{Mathematical Statistics and Probability} \\
	\vspace{.05in} 
	- Have knowledge of Bayesian statistics, point estimation, hypothesis testing, stochastic processes and martingale theory. 
	
	
	\item \textbf{Spanish} \\
	\vspace{.05in} 
	- Upper-intermediate level: speaking and writing.
	
\end{itemize}

\section{Certifications}
\begin{itemize}
	\item \textbf{Actuary Exams} \vspace{.05in}\\
	- P (Probability), FM (Financial Math) and MFE (Models for Financial Economics)
\end{itemize}

\section{Teaching Experience}
		* Instructor: Lectured, created and graded all exams \\  \vspace{.05in}
* GL = Only international students \\
* Recitation: Professor's assistant, gave weekly lectures, mentored students by appointment
\begin{itemize}[itemsep = -0pt, topsep=0pt]
	\item \textbf{Auburn University Graduate Teaching Assistant} {\hfill Aug 2021 - Aug 2022} \\ \vspace{.05in}

	- Transferred to Auburn University by the invitation of my advisor, Dr. Le Chen \\  \vspace{.05in}

	- Below is a list of my teaching assignments from AU \\  \vspace{.05in}
%	* Instructor: Lectured, created and graded all exams \\  \vspace{.05in}
%	* GL = only international students
		\begin{itemize}
			\item Math 1120-GL : Pre-Calculus I  (Instructor -in person - one section) \hfill Spring 2022 \\ \vspace{.05in}

			\item Head Tutor at Auburn University’s Math Tutoring Center \hfill Fall 2021 \\ \vspace{.05in}
			
			\item Stat 3600 : Probability and Statistics (Recitation - in person - one section)  \hfill Fall 2021 \\ \vspace{.05in}
%			 \hspace{.05in} - Supervised 14 tutors, created the weekly work schedule, assigned students to each tutor
		\end{itemize}
	
%		\begin{itemize}
%			\item In Fall 2021, I was the head tutor at Auburn University's math tutoring center. I was put in charge of the 14 other tutors and I would make the weekly schedule. I would also direct the incoming students to their assigned tutor and step in to tutor any subject that needed extra help for the day.
%			\item  In Spring 2022, I instructed a global course in Pre-Calculus. A global course means that all of the students in the class were from a different country. I would give lectures every Monday, Wednesday and Friday as well as hold weekly office hours. I also wrote and graded all of the exams.  
%		\end{itemize}
	
	\item \textbf{University of Nevada Las Vegas Graduate Assistant} {\hfill Aug 2017 - Aug 2021} \\ \vspace{.05in}

- Below is a list of my teaching assignments from UNLV \\ \vspace{.05in}

% * Instructor: Lectured, created and graded all exams \\  \vspace{.05in}
%
%* Recitation: Professor's assistant, gave weekly lectures, mentored students by appointment

		\begin{itemize}
			\item Math 124 : College Algebra  (Instructor - remote - one section) \hfill Summer 2021
			\item Math 458/658 : Intro to Real Analysis II (Recitation - remote - one section) \hfill Spring 2021
			\item Math 283 : Calculus III (Recitation - remote - one section) \hfill Spring 2021
			\item Math 181 : Calculus I (Recitation - remote - one section) \hfill Spring 2021
			\item Math 457/657 : Intro to Real Analysis I (Recitation - remote - two sections) \hfill Fall 2020
			\item Math 181 : Calculus I (Recitation - remote - two sections) \hfill Fall 2020
			\item Math 126 : Pre-Calculus I  (Instructor - remote - one section) \hfill Summer 2020
			\item Math 124 : College Algebra  (Instructor - in person - two sections) \hfill Spring 2020
			\item Math 124 : College Algebra  (Instructor - in person - two sections) \hfill Fall 2019
			\item UNLV Math Tutor (in person) \hfill Spring 2019
			\item Math 124 : College Algebra  (Instructor - in person - two sections) \hfill Fall 2018
			\item Math 126 : Pre-Calculus I  (Instructor - in person - one section) \hfill Summer 2018
			\item Math 126 : Pre-Calculus I  (Instructor - in person - two sections) \hfill Spring 2018
			\item UNLV Math Tutor (in person) \hfill Fall 2017
%			\item Instructed courses in College Algebra and Pre-Calculus 
%			\item Led recitation sections in Calculus I and III, and Real Analysis 
%			\item Tutored students from all levels in a group setting 
		\end{itemize}
\end{itemize}
  
\section{Acedemic Awards}
\begin{itemize}[itemsep = -0pt, topsep=0pt]
	    \item \noindent \textbf{Auburn University Research and Citation Award} \hfill Fall 2021 \\
	    \vspace{.05in} 
	- Presented by a committee of AU professors to a graduate student based on research.
	
    \item \noindent \textbf{Bhatnagar Math Award} \hfill Dec 2015 \\
    \vspace{.05in} 
     - A cash award presented by a committee of UNLV professors to a graduating undergraduate student based on math GPA. 
\end{itemize}


%-----------PROJECTS-----------------

 
 \section{Talks}
 
\begin{itemize}[itemsep = -0pt, topsep=0pt]
\item  \noindent Auburn University's Graduate Seminar \hfill April 27, 2022 \\
\vspace{.05in} 
- Presented an \href{https://github.com/nickeisenberg/School/blob/master/Presentations/GradSeminar(4-27-22)\%20/GdSemTalk.pdf}{"Introduction to Stochastic Analysis"} to first year mathematics \\ graduate students with the focus of introducing stochastic calculus.

\item \noindent Special session on Stochastic Analysis and Applications at the AMS Fall 2021 \\ Southeastern Sectional meeting \hfill November 20, 2021 \\
\vspace{.05in} 
- Invited to present some of my \href{https://github.com/nickeisenberg/School/blob/master/Presentations/AMSmeeting/ams.pdf}{thesis work}.
\vspace{.05in} 

\item \noindent Auburn University's Stochastic Seminar \hfill November 2, 2021 \\
\vspace{.05in} 
- Presented some of my \href{https://github.com/nickeisenberg/School/blob/master/Presentations/AUStochSeminar(11-20-21)/presentation.pdf}{thesis work} at the bimonthly Stochastic Seminar at AU. 
\end{itemize}

\section{Hobbies and Other Interests}
\begin{itemize}
	\item Rock Climbing - Indoor and outdoor bouldering and sport climbing. 
	\item Jazz Guitar - I really enjoy playing blues, bebop, bossa nova and funk. 
	\item Cooking 
\end{itemize}





%$\;$\vspace{.5in}
%
%
%\section{Personal Achievements}
%I would like to end the resume with some info about me and some of my achievements. 
%\begin{itemize}
%	\item Getting published. I started my PhD in 2018 and soon later I began working on several projects. I finished my first project in September of 2021 and my advisor and I choose to submit it to {\em Stochastics and Partial Differential Equation: Analysis and Computations} and it has just recently been accepted. Now I am currently working on the second bullet pointed project in my Publications section above. It has been quite challenging and I have hit many road blocks but I believe now that I am in a position where I may be able to submit very soon as I have recently been able to make significant progress. \vspace{.1in}
%	
%	\item Learning Spanish. To learn, I downloaded 150 30-minute Pimsleur tapes and listened to one or two of them each day. After that I started listening to podcasts in Spanish and watching Spanish television but with Spanish subtitles turned on. After about a year of doing this, I was confidently able to speak and have conversations with others in Spanish. 
%	
%	\item Learning jazz guitar. I started playing guitar when I was 19 and when I was 21, I took classical lessons for about a year and a half. I then remained self taught until I started taking jazz lessons in the summer of 2020. At that time, I started listening to jazz and I realized that this is the type of music I want to play. Jazz is the polar opposite to classical. It is highly creative and  almost 100\%  improvised. It is also insanely complex as you need to know how to improvise over numerous chord changes and to do so requires an intense amount of music theory and practice. \href{https://youtube.com/shorts/v5t6uV5slHM?feature=share}{Heres a clip of me improvising on a Bb jazz-blues.} \vspace{.1in}
%\end{itemize}
% 
% I really hope you consider me as a candidate for this quantitative-research position. I am well accustomed to research and I have published in a top journal in the field of stochastic partial differential equations. I am very easy to get along with, a very motivated worker and a very fast learner. I have been working on improving my programming skills by simulating various stochastic processes in Python, which can be seen on the GitHub link above, and I also have started studying Deep Learning. I am confident that I will be able to keep up with whatever is at task and I am passionate enough about this material to know that I will enjoy all the time I need to spend in order to be a valuable team-member at Point72. 
% 



%-------------------------------------------
\end{document}
